\documentclass[10pt,a4paper]{amsart}
\usepackage[margin=19mm]{geometry}
\usepackage{amsmath,amssymb,mathtools}
\usepackage[scr=rsfs,cal=euler]{mathalfa}
\usepackage{xfrac}
\usepackage[unicode,hidelinks,bookmarks=false]{hyperref}
\newcommand{\ie}{i.e.\ }
\newcommand{\cf}{cf.\ }
\newcommand{\resp}{respectively\ }
\newcommand{\Subsection}{\subsection}
\providecommand{\nobreakdash}{\nobreak\hbox{-}}
\setlength{\emergencystretch}{2em}
\multlinegap0pt
\usepackage{bm}
\usepackage{bbm}
\usepackage{dsfont}
\usepackage{xspace}
\usepackage{cancel}
\usepackage{mathtools}
\usepackage{amsthm}
\makeatletter
\@ifundefined{lemma}{\newtheorem{lemma}{Lemma}}{}
\@ifundefined{theorem}{\newtheorem{theorem}{Theorem}}{}
\@ifundefined{thmout}{\newtheorem{thmout}{Theorem}}{}
\makeatother
\def\D{\mathbb{D}}
\title[A note on forward iteration of inner functions]{A note on forward iteration of inner functions}
\author{Gustavo Rodrigues Ferreira}
\date{Source: arXiv:2206.00374v1 (CC BY 4.0)}
\begin{document}
\maketitle

\noindent\emph{Source and licence.} A note on forward iteration of inner functions. Version arXiv:2206.00374v1 (CC BY 4.0), distributed under the Creative Commons Attribution 4.0 International licence, as recorded in the arXiv licence metadata for this exact version and shown on the abstract page https://arxiv.org/abs/2206.00374v1. Full rights notice: licenses/arxiv-2206.00374-CC-BY-4.0.txt.\par\smallskip

\noindent\textit{Editorial scope.} Counted unit: the Lemma whose source label is \texttt{lem:estimates} in the article (source0.tex lines 147--151, source label \texttt{lem:estimates}), delivered together with the article's own complete original proof of it (source0.tex lines 152--186, one \texttt{proof} environment of 35 lines). No mathematics is written, repaired or re-derived: statement and proof are the article's own text line for line. Required context retained uncounted: thmout:class (lines 95--100); thm:inner (lines 107--110). Every definition, display and label of those lines is retained unchanged. Omitted from this package and not redistributed: 1--94, 101--106, 111--146, 187--428. Nothing inside a retained excerpt is omitted or altered: every retained body is delivered exactly as the article's lines are written, so no omission receipt of any kind accompanies this package. Citations: the retained excerpts carry no citation command at all, so no bibliography entry of the article is transcribed and no reference is redistributed. Reference adaptation: none was needed and none was made; every reference inside the delivered unit resolves inside it, so no unresolved reference or citation is produced. Printed slips kept literal: no printed word, symbol, hypothesis or number of the counted statement or of its proof was altered. Layout and class: the article's own class is \texttt{article} with packages including inputenc, enumerate, fontenc, hyperref, url, booktabs, amsfonts, amsmath, amssymb, amsthm, nicefrac, microtype, lipsum, xcolor; the wrapper is the lane's pinned \texttt{amsart} editorial wrapper at 10pt on A4, which restates the theorem-like environments the retained text needs and adds the article's own macro definitions that the retained text needs (\texttt{\textbackslash D}), with extra wrapper packages (bm, bbm, dsfont, xspace, cancel, mathtools). The delivered numbering is the unit's own. Assets: no retained span contains a figure environment, a graphics-inclusion command, a PDF-page inclusion, a table or a TikZ picture; no image file is copied into this package, the package declares no asset dependency and no image file is redistributed with it. Licence: version arXiv:2206.00374v1 of this article is distributed under the Creative Commons Attribution 4.0 International licence, as recorded in the arXiv licence metadata for this exact version and shown on the abstract page https://arxiv.org/abs/2206.00374v1. A full rights notice is delivered at licenses/arxiv-2206.00374-CC-BY-4.0.txt. Characters: every retained excerpt is pure ASCII, so the delivered engine, XeLaTeX with its default Latin Modern text font, reports no missing character.
\begin{thmout} \label{thmout:class}
Let $f_n:\D\to\D$ be holomorphic and such that $f_n(0) = 0$ for all $n\in\mathbb{N}$, and assume that the forward iteration of $f_n$ converges locally uniformly to a limit $f:\D\to\D$. Then, $f$ is non-constant if and only if
\begin{equation*}
    \sum_{n\geq 1} (1 - |f_n'(0)|) < +\infty.
\end{equation*}
\end{thmout}

\begin{theorem} \label{thm:inner}
Let $f_n:\D\to\D$, $n\in\mathbb{N}$, be inner functions such that $f_n(0) = 0$ for all $n\in\mathbb{N}$, and assume that $F_n := f_n\circ\cdots\circ f_1$ converges locally uniformly to some $F:\D\to\D$. Then, $F$ is either constant or an inner function, and the latter happens if and only if
\[ \sum_{n\geq 1} (1 - |f_n'(0)|) < +\infty. \]
\end{theorem}

\begin{lemma} \label{lem:estimates}
Let $g:\D\to\D$ be a holomorphic function with $g(0) = 0$ and $|g'(0)| = \lambda$. Then, for $z\in\D$ such that $|z| \leq \lambda$,
\[ |g(z)| \geq |z|\left(1 - \mu\cdot\frac{1 + |z|}{1 - |z|}\right), \]
where $\mu = 1 - \lambda$.
\end{lemma}

\begin{proof}[Proof of Theorem \ref{thm:inner}]
Let us assume that $F = \lim_{n\to +\infty} F_n$ is non-constant. It is clear from Weiestrass's theorem that $F$ is holomorphic; we must prove the stronger fact that it is an inner function. By Theorem \ref{thmout:class}, we have
\[ \sum_{n\geq 1} (1 - |f_n'(0)|) < +\infty; \]
in other words, $\prod_{n\geq 1} |f_n'(0)|$ converges to a positive number, and so given $\epsilon > 0$ there exists $N = N(\epsilon)$ such that
\[ \nu := \prod_{n=N+1}^{+\infty} |f_n'(0)| > 1 - \epsilon. \]
Now, for $n > N$, let
\[ H_n(z) := f_n\circ\cdots\circ f_{N+1}; \]
by Montel's theorem, there exists a subsequence $(H_{n_k})_{k\in\mathbb{N}}$ converging locally uniformly to some $H:\D\to\D$, and by Schwarz's lemma we in fact have $|H_n|\to |H|$. It follows that $|H'(0)| = \nu$, so that by Lemma \ref{lem:estimates} we have
\[ |H(z)| \geq |z|\left(1 - (1 - \nu)\frac{1 + |z|}{1 - |z|}\right) \geq |z|\left(1 - \epsilon\cdot\frac{1 + |z|}{1 - |z|}\right) \]
for $|z| < \nu$. Applying this to $|z| = 1 - \epsilon^{1/2} < 1 - \epsilon < \nu$, we get
\[ |H(z)| \geq (1 - \epsilon^{1/2})\left(1 - \epsilon\cdot\frac{2 - \epsilon^{1/2}}{\epsilon^{1/2}}\right) = (1 - \epsilon^{1/2})(1 - 2\epsilon^{1/2} + \epsilon), \]
and since $\epsilon > 0$ this gives
\begin{equation} \label{eq:H}
    |H(z)| \geq 1 - 3\epsilon^{1/2}.
\end{equation}
Define now the set $\Gamma_\epsilon := \{z : |F_N(z)| = 1 - \epsilon^{1/2}\}$; then, by (\ref{eq:H}),
\begin{equation} \label{eq:HGN}
    \text{$|F(z)| = |H\circ F_N(z)| \geq 1 - 3\epsilon^{1/2}$ for $z\in \Gamma_\epsilon$.}
\end{equation}
Furthermore, by Schwarz's lemma,
\begin{equation} \label{eq:Ge}
    \Gamma_\epsilon \subset \{z : 1 - \epsilon^{1/2} < |z| < 1\}.
\end{equation}
It is not true in general that for any inner function $f:\D\to\D$ fixing the origin the set $\{z : |f(z)| = r\}$ surrounds the circle $\{z : |z| = r\}$, even for $r$ very close to one; think, for instance, of the inner function $z\mapsto z\cdot\exp\left((z - 1)/(z + 1)\right)$. Nevertheless, we can show that $\Theta_\epsilon := \{\arg z : z\in\Gamma_\epsilon\}$ always has full measure in $[0, 2\pi)$, and that $\epsilon' < \epsilon$ implies $\Theta_\epsilon'\subset \Theta_\epsilon$.

Indeed, since $F_N$ is a composition of finitely many inner functions, it is itself an inner function, so that the set
\[ I_N := \{\theta\in[0, 2\pi) : \lim_{r\nearrow 1} |F_N(re^{i\theta})| = 1\} \]
has full Lebesgue measure. Applying the intermediate value theorem to the function $r\mapsto |F_N(re^{i\theta})|$, we see that for $\theta\in I_N$ there is always some $r^*\in (0, 1)$ such that $|F_N(r^*e^{i\theta})| = 1 - \epsilon^{1/2}$, and hence $I_N\subset \Theta_\epsilon$. Since $I_N$ has full measure in $[0, 2\pi)$, so does $\Theta_\epsilon$, as claimed. Now, if $\epsilon' < \epsilon$, we can assume that $N' = N(\epsilon') \geq N = N(\epsilon)$, so that by Schwarz's lemma $|F_{N'}(z)| \leq |F_N(z)|$ for all $z\in\mathbb{D}$. For $z' = r'e^{i\theta'}\in\Gamma_{\epsilon'}$, this means that $|F_N(z')| \geq |F_{N'}(z')| = 1 - \epsilon^{1/2}$, meaning that we can once again apply the intermediate value theorem and conclude that $\theta'\in \Theta_\epsilon$, and therefore $\Theta_{\epsilon'}\subset \Theta_\epsilon$.

Taking now a sequence $\epsilon_k \searrow 0$, we see that the set
\[ \Theta := \bigcap_{k\geq 1} \Theta_{\epsilon_k} \]
has full Lebesgue measure in $[0, 2\pi)$, and for $\theta\in\Theta$ we have by (\ref{eq:HGN}) and (\ref{eq:Ge}) that
\[ \limsup_{r\nearrow 1} |F(re^{i\theta})| = 1. \]
By Fatou's theorem, $F$ has well-defined radial limits in a full-measure set $E\subset [0, 2\pi)$, and the set $E\cap\Theta$ has full Lebesgue measure. For $\theta\in E\cap\Theta$, we must therefore have $|F(e^{i\theta})| = 1$; it follows that $F$ is an inner function.
\end{proof}

\end{document}
